\begin{textsample}{POS Dim 4 – pb – Score 7.00 – pb/pb\_em\_80.txt}  \label{ex:f4_pos_268}
2.4.5.2 \textbf{Territórios} e Fronteiras A \textbf{territorialidade} abrange muito mais do que uma \textbf{fronteira} física, ela se expande para além dos muros de concreto erguidos pelo preconceito e pelas idiossincrasias individuais e culturais.
Já as \textbf{fronteiras} culturais, morais, afetivas, emocionais e étnicas, segundo Santos ( 1996 ), chegam muito antes das \textbf{fronteiras} físicas.
Aqui, os fenômenos da globalização, da homogeneização e da desterritorialização são analisados do global para o local e do local para o global de forma dialética.
Atualmente, \textbf{territórios} e \textbf{fronteiras} passam a ter novas perspectivas diante da virtualização das relações sociais, das relações de poder, da intersubjetividade e da perspectiva da sociedade em rede digital, bem como da inteligência coletiva proporcionada pelo ciberespaço e pela cibercultura. 2.4.5.3 Natureza, Sociedade e Indivíduo Natureza talvez seja uma das categorias mais antigas e peculiares à área de Ciências Humanas e Sociais Aplicadas.
O conceito de natureza é um conceito histórico, geográfico, sociológico e filosófico.
Não há como tratar da questão da natureza sem perpassar por esses componentes.
Até porque, ao mesmo tempo em que estamos na natureza, somos natureza, como enfatiza Robert Lenoble ( 1990 ).
Desse modo, não há como falar ou definir esse conceito na perspectiva de fora, mas sempre de dentro.
Assim, garantir as liberdades individuais e os direitos fundamentais estabelecidos nos instrumentos normativos é a principal garantia de cidadania diante da diversidade cultural que nosso país abrange.
Nesse sentido, o conceito de sociedade, como uma categoria científica de análise da coletividade humana, passa pelo crivo da noção de natureza para ser contemplada no indivíduo como forma subjetiva de contemplação social.
Nesse sentido, a visão de natureza a ser vislumbrada vem em busca de um biocentrismo, ou seja, na construção de uma relação de respeito com a natureza que deve abranger todos os seres vivos no planeta. 2.4.5.4 Cultura, \textbf{Identidade} e Ética A cultura de um povo está estritamente ligada à nossa origem vivencial e à noção de cuidar.
O termo vem do latim “ colere ”, que significa cuidar.
Mais do que cuidar, cultura é um conjunto de convenções de normas, regras e valores que definem a personalidade de um povo.
O conceito de cultura como saber de um povo, faz do nosso \textbf{estado} um lugar mestiço que nos define como democráticos com múltiplas \textbf{identidades} em nossos comportamentos e vivências.
O processo de construção da \textbf{identidade} como definição de nossa singularidade e de nossa coletividade se encontra para além do eu e reflete no outro como alteridade moral e emocional.
A \textbf{identidade} cultural no nosso país é diversa, a julgar pelas diferenças étnicas e culturais que nos moldam como um povo multicultural.
Também, as diversidades ( étnica, de gênero, \textbf{geracional}, cultural, \textbf{religiosa} ).
Desse modo, o Brasil é um país diverso e multicultural que engloba várias perspectivas de sociabilidade, singularidade e culturalidade.
Abordar e evidenciar as diferentes culturas, manifestações artísticas e \textbf{religiosas}, visões de mundo e \textbf{territorialidades} existentes no país e no \textbf{Estado} da Paraíba é fundamental para que possamos contemplar, de maneira democrática e intercultural, as várias vivências que coexistem em nosso \textbf{estado}, país e planeta.

% matched lemmas: estado, fronteira, geracional, identidade, religioso, territorialidade, território
\end{textsample}
